% page 504 of the facsimile (image p-503 of the scan)
% block 157.5 x 237.7 mm ; left margin 8.0 mm
\pgc{-12.700mm}{0.000mm}{
\l{\cel{3}496.}
\l{}
\l{\sou{retretar}. - (\sou{netrans., de,}\cel{1}ek). I. (\sou{strategio)}. Retro-}
\l{\cel{3}iro (+\sou{rekulo)}\cel{2}da trupo de militisti qua koaktesas agar}
\l{\cel{3}tale da la enemiko plu forta qua avancas. - II. Cesar}
\l{\cel{3}sua vivo profesionala, stat-employatala, mestierala, mon-}
\l{\cel{3}dumala. - DEFIRS.}
\l{}
\l{\sou{retr}o- . Prefixo qua indikas movo de\cel{1}\sou{avan}\cel{1}ad\cel{1}\sou{dop}\cel{1}- de ula}
\l{\cel{3}epoko ad epoko antea. - L.}
\l{}
\l{\sou{retush}ar. (\sou{trans.}) Ri-laboreskar (ulo) por korektigar o plu-}
\l{\cel{3}bonigar ta o ca ek lua parti. - DEFIRS.}
\l{}
\l{\sou{reumatismo}. (\sou{patol.}) Dolorigo qua su manifestas sen inflamuro,}
\l{\cel{3}en la muskuli, en la artiki, e di qua la kauzo ne esas}
\l{\cel{3}determinita bone. - DEFIRS.}
\l{}
\l{\sou{revar}. (\sou{netrans., pri)}. Lasar sua penso vagar hazarde.- EF.}
\l{}
\l{\sou{revanchar}. (\sou{netrans., pri}). Retroprenar de ulu la avantajo}
\l{\cel{3}quan lu prenis de ni sen nia konsento. - DEFI.}
\l{}
\l{revelar. (trans., per, pri, ad).\cel{2}I. Konocigar (to quo esis}
\l{\cel{3}celita). - II. (\sou{teol.}) (\sou{aludante Deo}).Konocigar a la homi}
\l{\cel{3}per inspiro supernatura. - III. (\sou{fotogr.}) Eventigar reakto}
\l{\cel{3}kemiala qua igas la imajo aparar (sur la vitro-plako o sur}
\l{\cel{3}la filmo fotografala.) EFIS.}
\l{}
\l{\sou{revendikar}. (\sou{trans.}) Postular la retrodono, kom apartenajo}
\l{\cel{3}nia, di to quon kunhomo ditenas. - DFIS.}
\l{}
\l{\sou{revenuo}. (\sou{financ}o).To quon ulu obtenas, singla-yare, de}
\l{\cel{3}sua kapitalo, de sua domeno, de pensiono, de rento.-}
\l{\cel{3}DEF.}
\l{}
\l{\sou{reverberar}. (\sou{trans.}) (\sou{fiziko}).(\sou{aludante surfaco)}. Retro-}
\l{\cel{3}sendar la lumo-radii o la kaloro-radii qui frapas lu. -}
\l{\cel{3}DEFIS.}
\l{}
\l{\sou{reverencar}. (\sou{trans.)}\cel{2}Movar sua korpo por salutar ye ma-}
\l{\cel{3}niero qua manifestas respekto profunda. - DEFIRS.}
\l{}
\l{\sou{reverso.}\cel{2}I. La facio, opozata a la facio quan, en kozo,}
\l{\cel{3}onu regardas, o regardigas, prefere. - II. La facio di}
\l{\cel{3}medalio, di moneto-peco, opozata a la facio sur qua}
\l{\cel{3}esas reprezentita figuro. - DEFIS.}
\l{}
\l{\sou{reversionar}. (\sou{yuro-cienco}).\cel{2}Dicesas pri havaji qui retro-}
\l{\cel{3}iras a la persono qua donacabis li a kunhomo, kaze ke ica}
\l{\cel{3}mortas sen filii. - EFIRS.}
\l{}
\l{\sou{revizar}. (\sou{trans.}) Submisar itere (texto, proceso, kodexo}
\l{\cel{3}konto) ad exploro por reformar, korektigar, se necesa.}
\l{\cel{3}- DEFIRS.}
}
